% section: 離散と連続を結ぶ

\section{離散と連続を結ぶ}

\subsection{時間の幅を明示する}

微分方程式
\[
  \frac{dx}{dt}=F(x)
\]
を考える.
時間を幅 $h$ で区切り,
\[
  t_n=nh,\qquad a_n=x(t_n)
\]
と置く.
微分の定義から,
\[
  \frac{x(t_{n+1})-x(t_n)}h
  \approx \frac{dx}{dt}(t_n)
  =F(x(t_n))
\]
である.
従って,
\[
  a_{n+1}=a_n+hF(a_n)
\]
という漸化式が得られる.

$h$ は単なる記号ではない.
それは,一回の更新がどれくらいの時間を表しているかを決める.
もし $h=1$ と暗黙に置くなら,そのことを意識しておく必要がある.

\subsection{オイラー法の誤差を一歩だけ調べる}

テイラー展開により,
\[
  x(t+h)
  =x(t)+hx'(t)+\frac{h^2}{2}x''(t)+\cdots
\]
である.
オイラー法は,
\[
  x(t+h)\approx x(t)+hx'(t)
\]
と,二次以上の項を捨てている.
従って,一回の更新で生じる誤差は,通常 $h^2$ 程度である.

しかし,時間 $T$ まで進むには,およそ $T/h$ 回の更新が必要である.
一回あたり $h^2$ の誤差が累積すると,全体の誤差はおおよそ
\[
  \frac{T}{h}\cdot h^2=Th
\]
程度になる.
これが,オイラー法の全体誤差が $h$ の一次程度になる理由である.

厳密な誤差評価には条件が必要だが,細かい刻みを使うほど近似がよくなる理由はこのように説明できる.

\subsection{指数関数は連続な等比数列である}

最も基本的な微分方程式
\[
  \frac{dx}{dt}=\lambda x
\]
を考える.
オイラー法では,
\[
  a_{n+1}=a_n+h\lambda a_n
  =(1+h\lambda)a_n
\]
となる.
これは等比数列である.
\[
  a_n=(1+h\lambda)^n a_0
\]

一方,微分方程式の解は,
\[
  x(t)=e^{\lambda t}x(0)
\]
である.
$t=nh$ とすれば,
\[
  e^{\lambda nh}=\left(e^{\lambda h}\right)^n
\]
だから,連続的な指数関数も,細かい時間幅で見れば等比的な更新の極限として現れる.

実際,
\[
  e^{\lambda h}
  =1+\lambda h+\frac{\lambda^2h^2}{2}+\cdots
\]
なので,$h$ が小さいとき,
\[
  1+\lambda h
\]
は $e^{\lambda h}$ の一次近似である.

「等比数列」と「指数関数」は別々の公式ではない.
一段ずつ進む世界と,連続的に進む世界で,同じ増加の仕組みを記述している.

\subsection{離散的な安定性と刻み幅}

$\lambda<0$ のとき,連続解 $e^{\lambda t}$ は $0$ へ近づく.
オイラー法の漸化式
\[
  a_{n+1}=(1+h\lambda)a_n
\]
が $0$ へ近づくためには,
\[
  |1+h\lambda|<1
\]
が必要である.
従って,
\[
  0<h<-\frac2\lambda
\]
でなければならない.

元の微分方程式が安定でも,刻み幅を大きく取りすぎると,近似の漸化式は発散する.
これは,「連続化したから安心」ではないことを示している.
離散化は新しい漸化式を作る操作であり,その漸化式自身の性質を調べなければならない.
